\section{Related Work}

\subsection{Benchmarks for Image Editing}

% \begin{table}[t]
% \caption{\textbf{Comparison between \bench\ and existing image editing benchmarks.} 
% \texttt{AVR} denotes Algorithm Visual Reasoning, \texttt{MIA} denotes Multi-Image Awareness, \texttt{WKR} denotes World Knowledge Reasoning, \texttt{DM} denotes Dynamic Manipulation, and \texttt{HP} denotes Human Preference evaluation.}
% % \vspace{0.3em}
% \label{tab:compare_image_editing_bench}
% \centering
% \resizebox{0.80 \linewidth}{!}{%
% \begin{tabular}{lcccccccc}
% \toprule
% \textbf{Benchmark} & \textbf{Size} & \textbf{Tasks} & \textbf{Lang.} & \textbf{AVR} & \textbf{MIA} & \textbf{WKR} & \textbf{DM}& \textbf{HP} \\
% \midrule

% Magicbrush~\cite{zhang2023magicbrush} & 1,053 & 5 & EN & \icono & \icono &\icono & \icono &  Low \\
% Emu Edit~\cite{cui2025emu3} & 3,055 & 7 & EN & \icono & \icono & \icono & \icono  & Low \\
% AnyEdit-Bench~\cite{yu2025anyedit} & 1,250 & 25 & EN & \icono & \icono & \icono & \icono & Low \\
% ImgEdit-Bench~\cite{ye2025imgedit} & 811 & 11 & EN & \icono & \icono & \icono & \icono & Mid \\
% GEdit-Bench~\cite{liu2025step1x} & 606 & 14 & EN/ZH & \icono & \icono & \icono & \icono & Mid \\
% ICE-Bench~\cite{pan2025ice} & 6,538 & 31 & EN & \icono & \icono & \icono & \icono & Low \\
% RISE-Bench\cite{zhao2025envisioning} & 360 & 4 & EN & \icono & \icono & \icohalf & \icono  & Mid \\
% UniREditBench~\cite{han2025unireditbench} & 2,700 & 13 & EN & \icono & \icono & \icono & \icohalf  & Mid \\
% WiseEdit~\cite{pan2025wiseedit} & 1,220 & 13 & EN/ZH & \icono & \icohalf  & \icoyes & \icono  & Mid \\
% \midrule
% \bench~(Ours) & 2,388 & 36 & EN/ZH & \icoyes & \icoyes & \icoyes & \icoyes  & High \\
% \bottomrule
% \end{tabular}%
% }
% % \vspace{-1.70em}
% \end{table}

\begin{table}[t]
 
\caption{\textbf{Comparison between \rmbench\ and existing image editing reward benchmarks.}
\texttt{AVR} denotes Algorithm Visual Reasoning, \texttt{DM} denotes Dynamic Manipulation, \texttt{WKR} denotes World Knowledge Reasoning, and \texttt{CP} denotes Complex Paint.}
\vspace{0.5em}
\label{tab:compare_image_editing_reward_bench}
\centering
\resizebox{0.95\linewidth}{!}{%
\begin{tabular}{lccccccccc}
\toprule
\textbf{Benchmark} & \textbf{Size} & \textbf{Tasks} & \textbf{Sampling Strategy}  & \textbf{AVR} & \textbf{DM} & \textbf{WK} & \textbf{CP} \\
\midrule
GenAI-Bench~\cite{li2024genai} & 919 & 7 & Cross-model  & \icono & \icono & \icono & \icono  \\
EditReward-Bench~\cite{luo2025editscore} & 3,072 & 13 & Cross-model & \icono & \icono & \icono & \icono  \\
EditReward-Bench~\cite{wu2025editreward} & 1,500 & 8 & Cross-model  & \icono & \icono & \icono & \icono  \\
MMBench2~\cite{hu2025multimodal} & 1,000 & - & Cross-model   & \icono & \icono & \icohalf & \icono  \\
\midrule
\rmbench~(Ours) & 2,251 & 36 & Cross-/Intra-model   & \icoyes & \icoyes & \icoyes & \icoyes  \\

\bottomrule
\end{tabular}%
}

\end{table}

Existing image editing benchmarks face two major limitations: limited task coverage and insufficient evaluation reliability. As shown in Table~\ref{tab:compare_image_editing_bench}, early benchmarks~\cite{zhang2023magicbrush,sheynin2024emu,yu2025anyedit,pan2025ice} mainly focus on narrow editing tasks and rely on automated metrics such as CLIP-I and DINO-I. However, these metrics often fail to capture fine-grained editing quality, especially for tasks involving world knowledge, visual consistency, and complex instruction following.
Recent benchmarks~\cite{ye2025imgedit,liu2025step1x,zhao2025envisioning,zhang2026well} adopt powerful MLLMs as judges for more flexible evaluation. Nevertheless, their reliance on simple judging prompts can lead to unstable assessments and misalignment with human judgment in complex scenarios.
To address these limitations, we propose \textbf{{\bench}}, a comprehensive benchmark covering $36$ fine-grained tasks across six categories. Beyond broad task coverage, {\bench} introduces human-aligned evaluation prompts with structured reasoning and carefully designed scoring rubrics, enabling more accurate, reliable, and interpretable assessment of image editing models.


\subsection{Benchmarks for Image Editing Reward Model}



With the rapid progress of image generation and editing, reward models have become increasingly important for improving instruction following and visual consistency through reinforcement learning (RL). Accordingly, reliable evaluation of image editing reward models has attracted growing attention.
As shown in Table~\ref{tab:compare_image_editing_reward_bench}, existing benchmarks~\cite{luo2025editscore,wu2025editreward,hu2025multimodal} typically construct preference pairs from limited editing tasks or from outputs generated by different models. However, such settings often deviate from practical RL scenarios, where reward models are required to compare candidate outputs produced by the same editing model under the same instruction. This mismatch limits faithful assessment of reward model quality and training effectiveness.
Recent efforts~\cite{zhao2025envisioning,deng2025emerging} have expanded evaluation coverage to more diverse tasks, including world knowledge and visual reasoning. Nevertheless, existing benchmarks still lack realistic and controlled preference construction, particularly in balancing task diversity and comparison consistency.
To bridge this gap, we propose \textbf{{\rmbench}}, a comprehensive benchmark for evaluating image editing reward models. {\rmbench} constructs preference pairs under more realistic and controlled settings, enabling multidimensional analysis of reward models in terms of instruction following, visual consistency, perceptual quality, and reasoning-aware editing preference.


% \begin{table}[t]
 
% \caption{\textbf{Comparison between \rmbench\ and existing image editing reward benchmarks.}
% \texttt{AVR} denotes Algorithm Visual Reasoning, \texttt{DM} denotes Dynamic Manipulation, \texttt{WKR} denotes World Knowledge Reasoning, and \texttt{CP} denotes Complex Paint.}
% \vspace{0.5em}
% \label{tab:compare_image_editing_reward_bench}
% \centering
% \resizebox{0.65\linewidth}{!}{%
% \begin{tabular}{lccccccccc}
% \toprule
% \textbf{Benchmark} & \textbf{Size} & \textbf{Tasks} & \textbf{Sampling Strategy}  & \textbf{AVR} & \textbf{DM} & \textbf{WK} & \textbf{CP} \\
% \midrule
% GenAI-Bench~\cite{li2024genai} & 919 & 7 & Cross-model  & \icono & \icono & \icono & \icono  \\
% EditReward-Bench~\cite{luo2025editscore} & 3,072 & 13 & Cross-model & \icono & \icono & \icono & \icono  \\
% EditReward-Bench~\cite{wu2025editreward} & 1,500 & 8 & Cross-model  & \icono & \icono & \icono & \icono  \\
% MMBench2~\cite{hu2025multimodal} & 1,000 & - & Cross-model   & \icono & \icono & \icohalf & \icono  \\
% \midrule
% \rmbench~(Ours) & 2,251 & 36 & Cross-/Intra-model   & \icoyes & \icoyes & \icoyes & \icoyes  \\

% \bottomrule
% \end{tabular}%
% }
% \vspace{-1.5em}
% \end{table}